\begin{proof}
\emph{Exact first-column coordinates.}
Put $t=s^2$, with $s$ real near zero.  Every contraction $V$ in
\eqref{I:R:eq:hk-definition} can be written
\begin{equation}
 V=[sz,W_s(z)N],\quad \Rnorm z=1,\quad \Rnorm N\le1,\qquad
 W_s(z)=I-(1-\sqrt{1-s^2})zz^*.
 \label{I:R:eq:first-column-chart}
\end{equation}
Indeed, the Schur complement of $I-V^*V$ says that the remaining columns
$V_F$ satisfy
\[
 V_F^*\left(I+\frac{s^2}{1-s^2}zz^*\right)V_F\preceq I;
\]
its inverse positive square root is $W_s(z)$.  Partition
\[
 z=\begin{pmatrix}z_0\\z_F\\z_o\end{pmatrix},\qquad
 N=\begin{pmatrix}c\\D\\l\end{pmatrix},
 \qquad z_F\in\RC^m,
\]
where $c$ is a row.  Define $J(s,z,N)=\ell B^{m+1}e_1/s$ by its
continuous polynomial extension at $s=0$, and
$M(s)=\max_{z,N}|J(s,z,N)|$.  Then
\begin{equation}
 h_k(s^2)=s^2M(s)^2,\qquad J(0,z,N)=lD^mz_F,\qquad M(0)=\Gamma.
 \label{I:R:eq:M-and-J}
\end{equation}
At every leading maximizer, Lemma~\ref{I:R:lem:cyclic-lossless-maximizer}
forces $z_0=z_o=0$, $\Rnorm{z_F}=1$, $D^*D+l^*l=I_m$, and $c=0$.
Thus the leading maximizers are precisely the embedded scalar boundary
set, up to an output phase.

\emph{Uniform localization of the normal directions.}
Expansion of the matrix product gives
\begin{align}
 J(s,z,N)&=lD^mz_F+sJ_1(z,N)+O(s^2),
 \label{I:R:eq:J-first-order}\\
 J_1(z,N)&=z_ocD^{m-1}z_F+z_0lD^{m-1}z_F
       +\sum_{j=0}^{m-2}(lD^jz_F)(cD^{m-2-j}z_F).
 \label{I:R:eq:J1-explicit}
\end{align}
The remainder is uniform on the compact parameter domain.  Put
$\delta=\Gamma^2-|lD^mz_F|^2$.  Near the leading maximizers,
\[
 |z_0|^2+|z_o|^2\le\frac\delta{\Gamma^2},\qquad
 \Rnorm c^2\le\operatorname{tr}(I-D^*D-l^*l)\le C\delta
\]
by Lemma~\ref{I:R:lem:quantitative-defect-control}.  Therefore
$|J_1|\le C\sqrt\delta$ there.  Away from a fixed neighborhood of the
leading maximizers, $\delta$ is bounded below, so the same estimate holds
globally after enlarging $C$.  It follows that
\begin{equation}
 |J(s,z,N)|\le\Gamma-\frac\delta{2\Gamma}
                         +C|s|\sqrt\delta+Cs^2.
 \label{I:R:eq:normal-gap-upper}
\end{equation}
A fixed leading maximizer gives $M(s)\ge\Gamma-Cs^2$.
Completing the square in \eqref{I:R:eq:normal-gap-upper} proves
\begin{equation}
 M(s)=\Gamma+O(s^2),\qquad
 \delta=O(s^2)\quad\hbox{at every global maximizer.}
 \label{I:R:eq:normal-localization}
\end{equation}
In particular, all maximizing $z_0,z_o,c$ are $O(s)$ and their remaining
passivity defect is $O(s^2)$.

\emph{Polar coordinates and the complex phase.}
At any maximizing point write
\[
 z_0=s\lambda,\quad z_o=s\zeta,\quad c=sw,\quad
 z_F=\sqrt{1-s^2(|\lambda|^2+|\zeta|^2)}\,x,\qquad \Rnorm x=1.
\]
Let $Y_{\rm raw}=\begin{pmatrix}D_{\rm raw}\\l_{\rm raw}\end{pmatrix}$
denote the lower block of $N$, and put
$T=I-Y_{\rm raw}^*Y_{\rm raw}$.  Then
$T\succeq c^*c$, $\operatorname{tr}T=O(s^2)$, and the polar
factorization is
\begin{equation}
 Y_{\rm raw}=U(I-T)^{1/2},\quad U^*U=I_m,\quad
 T=s^2(w^*w+F),\qquad F\succeq0.
 \label{I:R:eq:normal-polar}
\end{equation}
The variables $\lambda,\zeta,w,F$ are uniformly bounded at maximizing
points.  Write $U=\begin{pmatrix}D\\l\end{pmatrix}$.  The slow-coordinate
unitary $\operatorname{diag}(e^{i\psi},I_m)$ preserves feasibility and
rotates both $J$ and $lD^mx$ by $e^{i\psi}$.  It therefore permits the
normalization $a=lD^mx>0$.  More explicitly, it changes
$x$ to $e^{i\psi}x$, $c$ to $e^{-i\psi}c$, $z_o$ to
$e^{i\psi}z_o$, and leaves $z_0,D,l$ unchanged.  Thus this normalization
discards no complex direction.  The base triples $(D,l,x)$ approach
$\mathcal M_m$, and always satisfy $a\le\Gamma$.

\emph{The normal quadratic form.}
Expand $(I-T)^{1/2}$, $W_s(z)$, and the finite product $\ell B^{m+1}e_1/s$.
Uniformly over the localized parameters,
\begin{equation}
 |J(s)|=a+s^2Q+O(s^3).
 \label{I:R:eq:J-normal-expansion}
\end{equation}
At a boundary triple in $\mathcal M_m$, the quadratic coefficient is
\begin{align}
 Q={}&-\frac12\operatorname{Re}S
   +\operatorname{Re}(\lambda a_{m-1})-\frac\Gamma2|\lambda|^2
   +\operatorname{Re}(wv)-\frac12wKw^*\nonumber\\
   &\hspace{18mm}-\frac\Gamma2|\zeta|^2
                       -\frac12\operatorname{tr}(KF).
 \label{I:R:eq:normal-quadratic-Q}
\end{align}
Off the boundary set the same formula uses the continuously varying
phase-normalized coefficients and replaces $\Gamma$ in the two
normalization costs by $a$.

For completeness, each contribution in
\eqref{I:R:eq:normal-quadratic-Q} can be read directly from the block
product.  Normalizing $z_F$ contributes
$-\Gamma(|\lambda|^2+|\zeta|^2)/2$.  The slow diagonal entry
$s^2\lambda$ contributes $\lambda a_{m-1}$.  A single slow-to-fast return
through the row $sw$ contributes
\[
 \sum_{j=0}^{m-2}(lD^jx)(wD^{m-2-j}x)=wv.
\]
Right contraction by $(I-s^2(w^*w+F))^{1/2}$ contributes
\[
 -\frac12\operatorname{Re}\sum_{j=0}^{m}
 lD^j(w^*w+F)D^{m-j}x
 =-\frac12wKw^*-\frac12\operatorname{tr}(KF).
\]
The first-column Schur factor contributes the perturbation
$D\mapsto D-s^2xx^*D/2$, hence $-S/2$.  Finally the slow readout is
$s^2\zeta$; using it requires a return to the slow state and first affects
$J$ at order $s^3$.  These observations prove the displayed expansion,
including its empty-sum case $m=1$.

Lemma~\ref{I:R:lem:response-K-positive} allows the quadratic form to be
maximized explicitly:
\begin{equation}
 \begin{gathered}
 F=0,\quad\zeta=0,\quad
 \lambda=\frac{\overline{a_{m-1}}}{\Gamma},\quad w=v^*K^{-1},\\
 Q_{\max}=-\frac12\operatorname{Re}S
       +\frac{|a_{m-1}|^2}{2\Gamma}+\frac12v^*K^{-1}v.
 \end{gathered}
 \label{I:R:eq:normal-quadratic-max}
\end{equation}
For the global upper bound, every sequence of maximizers has a
subsequence whose base triples approach $\mathcal M_m$.  Their base
contribution $a-\Gamma$ is nonpositive, all quadratic coefficients
converge continuously, and $K$ remains uniformly positive.  Therefore
\[
 \limsup_{s\to0}\frac{M(s)-\Gamma}{s^2}
 \le\max_{\mathcal M_m}Q_{\max}.
\]
For the reverse bound choose a boundary triple maximizing
$Q_{\max}$, take the parameters in
\eqref{I:R:eq:normal-quadratic-max}, and use the exact square roots in
\eqref{I:R:eq:first-column-chart} and \eqref{I:R:eq:normal-polar}.
They form a feasible family for all sufficiently small $|s|$ and give
the reverse liminf.  Consequently
\[
 M(s)=\Gamma+s^2\max_{\mathcal M_m}Q_{\max}+o(s^2).
\]
Equation~\eqref{I:R:eq:M-and-J} proves the first expansion in
\eqref{I:R:eq:all-budget-envelope-expansion}, with the coefficient
\eqref{I:R:eq:bk-response-formula}.

\emph{The final scalar maximization.}
The elementary estimate $h_k(t)\le t$ and
\[
 1-\sqrt{(1-q)(1-t)}\ge\frac{q+t}{4}
\]
show that the objective in \eqref{I:R:eq:allk-elimination} is at most
$16qt/(q+t)^2$.  The trial point $t=q$ has value tending to
$\Gamma^2>0$, so all maximizers have $x=t/q$ in a fixed compact
subinterval of $(0,\infty)$ when $q$ is small.  Uniformly on such
intervals,
\begin{align*}
 \frac{q h_k(qx)}{(1-\sqrt{(1-q)(1-qx)})^2}
 ={}&\frac{4\Gamma^2x}{(1+x)^2}\\
 &+q\left(\frac{4b_kx^2}{(1+x)^2}
        -\frac{2\Gamma^2x(x-1)^2}{(1+x)^3}\right)+o(q).
\end{align*}
The leading term has its unique maximum $\Gamma^2$ at $x=1$.
Thus all maximizing $x$ tend to one.  The leading term is never above
$\Gamma^2$ and the bracket tends to $b_k$ along every maximizing
sequence.  Taking $x=1$ gives the matching lower bound, proving the
second expansion.  The established projection inequality
$D_k(q)\ge1-\mathcal P_k(q)$ proves
\eqref{I:R:eq:all-budget-D-lower}.
\end{proof}
